\documentclass[11pt]{article}
\usepackage{fontspec,amsmath,amssymb,amsthm,hyperref}
\usepackage[margin=1in]{geometry}
\setlength{\emergencystretch}{3em}
\title{The actual gradient generator and its derived L\ensuremath{{}^{2}} bound}
\author{LangevinLean mathematical exposition}
\date{}
\begin{document}\maketitle
\noindent Source module: \texttt{AverageLMC/LaplacianGenerator.lean}. Compile this proof document with XeLaTeX or LuaLaTeX. The accompanying HTML pairs the mathematical argument with the exact Lean statements and records the scope of the formalization.


The vector generator in the third-order Langevin argument is represented by true derivatives of the same potential. Its continuity, finite second moments and quantitative L\ensuremath{{}^{2}} bound follow from the paper’s C\ensuremath{{}^{3}}, gradient-Lipschitz and squared Laplacian-gradient assumptions.



The actual vector generator of the gradient and its sharp stationary \(L^2\) bound are


\[G_f(x)=\nabla\Delta f(x)-D(\nabla f)(x)[\nabla f(x)]=\Delta(\nabla f)(x)-D(\nabla f)(x)[\nabla f(x)].\]
\[Y\sim\pi\quad\Longrightarrow\quad\|G_f(Y)\|_{L^2}\le L\sqrt{M_{\rm av}p}+M_\Delta\sqrt p.\]
\section{Define and identify the true generator}


The vector field is ∇Δf−D(∇f)∇f. The directional-derivative trace formula makes Δf continuously differentiable for C\ensuremath{{}^{3}} f. The Riesz representation of its derivative gives continuity of ∇Δf; evaluation of the continuous Hessian on the continuous gradient gives continuity of the second term. The proved vector commutation identity identifies this field exactly with Δ(∇f)−D(∇f)∇f, the Langevin generator applied to ∇f.



\noindent Lean declarations: \texttt{gradientGenerator}, \texttt{contDiff\_laplacian\_one}, \texttt{continuous\_gradientGenerator}, \texttt{gradientGenerator\_eq}.

\section{Bound the two actual derivative terms pointwise}


A globally L-Lipschitz gradient has Fréchet derivative operator norm at most L at every point. Its value on ∇f is therefore bounded by L‖∇f‖. Separately, the paper’s hypothesis ‖∇Δf‖\ensuremath{{}^{2}}\ensuremath{\le}Mδ\ensuremath{{}^{2}}p implies ‖∇Δf‖\ensuremath{\le}Mδ√p because both norms and Mδ are nonnegative. Neither bound asks for the Hessian to be Lipschitz.



\noindent Lean declarations: \texttt{hessian\_gradient\_norm\_le}, \texttt{gradient\_laplacian\_norm\_le}.

\section{Construct finite second moments for the generator}


The Laplacian-gradient term is a continuous function bounded uniformly by Mδ√p, hence belongs to L\ensuremath{{}^{2}} on every probability space. The Hessian-gradient term is continuous and bounded by L times the actual gradient; the latter has finite second moment because it is a Lipschitz field applied to an L\ensuremath{{}^{2}} state. Subtraction constructs the finite second moment of the actual generator itself. These are conclusions of the potential assumptions, not moment premises supplied to a local-error theorem.



\noindent Lean declarations: \texttt{laplacian\_gradient\_memLp}, \texttt{hessian\_gradient\_memLp}, \texttt{gradientGenerator\_memLp}.

\section{Derive the sharp stationary generator constant}


Represent each generator term in L\ensuremath{{}^{2}} and use the exact almost-everywhere pointwise identities. The first representative has norm at most Mδ√p; the second has norm at most L times the lifted score norm. The triangle inequality therefore bounds the actual generator norm by L‖∇f(Y)‖\ensuremath{{}_{2}}+Mδ√p. When Y has the true Gibbs law, pushforward integration identifies its score second moment with the Gibbs integral. The sharp coordinate-score theorem bounds that integral by Mav p. Taking its nonnegative square root yields the required stationary constant L√(Mav p)+Mδ√p.



\noindent Lean declarations: \texttt{lp\_gradientGenerator\_norm\_le}, \texttt{lp\_gradientGenerator\_gibbs\_norm\_le}.

\end{document}
